\chapter{AI-Assisted Research and Verification Methodology}
\label{ch:ch15_methodology}

The question of this chapter is how the statements of this manuscript were produced and what kind of evidence stands behind each of them.
\Cref{sec:ai} describes the AI-assisted workflow, the author's role and the limits of AI review. \Cref{app:methods} describes the verification:
the classes of evidence, the scripts and what a passing check shows, and the parts that have no executable evidence. The verification
manifest, one row per check with its script, command, method and supporting label, is \cref{app:manifest}.

\section{The AI-assisted workflow}
\label{sec:ai}

\paragraph{How the work was done.} By the author's account, the programme was developed over roughly three years in an iterative
process with AI systems. The questions, the starting hypothesis and the direction of the programme are the author's. Claude models
assisted extensively: with candidate theorems, exploratory proofs, alternative constructions, code, calculations and drafting of the prose.
ChatGPT assisted with discussion, methodology and referee-style criticism; the five reviews and the temporal audit consolidated in the review
checklist that governs the present revision were AI-assisted reviews of this kind. The author examined the results throughout, raised
objections and asked for corrections; several of the corrected statements are recorded in \cref{app:methods}. The model versions changed over
the period, and no particular version was in use throughout. The repository history of the programme begins in June 2026.

\paragraph{Roles in the present revision.} In the revision that produced this manuscript, AI sessions under the author's direction:
\begin{enumerate}[label=(\roman*),leftmargin=2.2em]
\item wrote derivation scripts together with the sections they check;
\item wrote referee scripts in separate sessions, instructed to refute the claims, not to import the derivation code, and to use a different
method where possible (the surviving instructions are kept in the author's records);
\item adjudicated the items of the review checklist, each with a recorded outcome and evidence;
\item read and compared the literature of \cref{ch:ch14_related}, recording what was read where;
\item assembled and edited this manuscript.
\end{enumerate}
The independence of a referee script is therefore an independence of code and method, not of a second human checker. None of the check
scripts is known to have been written by a second human, and the manuscript does not claim that every result was checked by a human or by a
second implementation.

\paragraph{What AI review is not.} The reviews that shaped the revision were AI-assisted. They are not formal external peer review, they
certify nothing, and agreement between models, or between runs of one model, is not a proof. A statement of this manuscript is supported
by its own evidence, of one of the classes of \cref{app:methods}, and not by the fact that an AI system produced or accepted it. The
manuscript has not been peer reviewed.

\paragraph{Responsibility.} The author is responsible for the manuscript, its verification and its conclusions. The disclosure above is meant
to say how the work was done; it implies no endorsement, coauthorship or guarantee by any AI system or its provider.

\section{How the statements were verified}
\label{app:methods}
\label{s9:sec:methods}
\label{r9:sec:main}

\subsection{Classes of evidence}

Each statement belongs to one of the following classes, and the class is a property of the evidence, not of the confidence in it.
\begin{enumerate}[label=(\roman*),leftmargin=2.2em]
\item \emph{Written proof.} An argument given in the text, which a reader can check line by line. Many of these were first drafted with AI
assistance; they stand or fall with the argument as written.
\item \emph{Exact certificate.} An identity, rank, dimension, spectrum or decomposition computed in integer or rational arithmetic, or as a
polynomial identity in symbolic variables. Every generator of the realisation of \cref{app:realisation} is a signed permutation matrix, so
every blade at every rung is one, and products, commutators and traces involve no rounding. Kernel dimensions such as $272$, centraliser and
stabiliser dimensions and the orbit dimensions $34$ and $20$ are exact ranks.
\item \emph{Numerical check.} A floating-point computation with a stated tolerance: integrations of the interior and orbital dynamics,
eigenvalue computations, random tests at seeded inputs, and minimisations. It confirms an exact formula or supports a statement that is
labelled numerical; it proves nothing beyond the sampled inputs.
\item \emph{Independent re-check.} A referee script, written without importing the derivation script and where possible by a different
method. The two share the definitions of the realisation and the conventions of the text, so a convention error common to both would not be
caught; what they derive separately is the code and the method of computation.
\item \emph{Citation.} An established result of the literature, used with its hypotheses and not recomputed: anomaly cancellation of the
Standard Model, the Vafa--Witten theorem, Michel's theorem, the Helmholtz conditions, Lanczos's identity, O'Neill's formulas and the others
marked inherited.
\item \emph{AI review.} A report of an AI-assisted reviewer. It produced objections, which were adjudicated with evidence of the classes
above; it is not itself evidence.
\end{enumerate}

\subsection{The scripts}

Each script prints one line \texttt{PASS <id>: <claim>} or \texttt{FAIL <id>: <claim>} per check and exits with status zero only if every check
passes. They need Python~3 with NumPy, SciPy, SymPy and mpmath and run from any directory, for example
\begin{center}
\texttt{python3 verification/run\_all.py --out reproducibility/selected-results/v1.0.0 -j 3}
\end{center}
which runs all of them and writes one record per script. The scripts are kept by chapter in \path{verification/chNN_<topic>/}; by origin they are:
\begin{enumerate}[label=(\roman*),leftmargin=2.2em]
\item \emph{Two readings} (\cref{ch:ch03_readings}): derivation scripts \texttt{r1}--\texttt{r4} and referee scripts \texttt{s1}--\texttt{s4}. The
Grassmann-valued Dirac field is handled by an exact Grassmann jet algebra in \texttt{r1}, and independently in \texttt{s1} by polarisation, which
encodes a quadratic Grassmann expression as an antisymmetric bilinear in two commuting fields. Nine checks are floating-point: r2-10c, r2-10d
and the random part of r2-18; s1-15, s1-31, s1-38; s2-10b, the random part of s2-14b, and s2-15.
\item \emph{Rest energy} (\cref{ch:ch03_readings,ch:ch04_rigidbody,ch:ch05_tower,ch:ch09_bfield}): derivation scripts \texttt{d1}--\texttt{d7} and
referee scripts \texttt{x1}--\texttt{x7}. Spectra and Jordan structure are checked with rational data, since floating-point inputs give false
negatives in semisimplicity tests; dynamics is integrated with relative tolerances $10^{-11}$ to $10^{-12}$, and elliptic and Mathieu quantities
use \texttt{mpmath} and \texttt{scipy.special}. The referee findings of the second part, fifty in number, were each re-examined by a second
referee session: forty confirmed, ten in part, none rejected, and no theorem refuted.
\item \emph{Emergence} (\cref{ch:ch05_tower,ch:ch09_bfield,ch:ch10_classical}): derivation scripts \texttt{e1}--\texttt{e7}; referee scripts
\texttt{y1}, \texttt{y2}, \texttt{y5}, \texttt{y6} for four of the seven sections; the review adjudications \texttt{z1} and \texttt{z5}, which also hold the
exact rank certificate for LAD-02 in place of an earlier floating-point rank. Two sections were refereed by reading only. The classical section
was corrected and reviewed in two passes; its observational inputs were checked against the primary papers where these were reachable, and the
others were not rechecked.
\item \emph{Photon selection} (\cref{ch:ch08_photon}): derivation scripts \texttt{p1}--\texttt{p3} and the adjudication script \texttt{z2}, which
rebuilds the $512$-dimensional realisation independently. No referee script exists for \texttt{p1} or \texttt{p2}. The global minimisations
p2-16 and p3-22 are numerical.
\item \emph{Global manuscript} (\cref{ch:ch01_clifford,ch:ch06_chiral,ch:ch07_vacuum}): the scripts \texttt{z3} and \texttt{z4} (\path{verification/ch06_chiral_matter}, \path{verification/ch07_mass_vacuum}), written to adjudicate the CORE review items; \texttt{z4} computes the Molien--Weyl counts of invariants on the
kernel space in exact integer arithmetic.
\item \emph{Rigid body} (\cref{ch:ch04_rigidbody}): \path{verification/ch04_rigid_body/rb1}--\texttt{rb3}, including the numerical orbital examples with
their inputs, solver and tolerances.
\item \emph{Release audit} (\cref{ch:ch05_tower,ch:ch07_vacuum,ch:ch08_photon}): \texttt{a6}--\texttt{a8}, written for version 1.0.0 for the statements
that the audit of \cref{app:repro} had left without a recomputation. \texttt{a6} and \texttt{a7} are exact (blade parity, integer arithmetic on the
$512$-dimensional module, exact modular ranks, Molien--Weyl integrals); \texttt{a8} is numerical, on the kernel space of \cref{ch:ch07_vacuum}, and
asserts a spectral gap at every rank it reads.
\end{enumerate}
The independent review checks (\path{verification/ch15_ai_methodology}) add $22$ checks, the vacuum checks (\path{verification/ch07_mass_vacuum},
v1--v3) $31$, the gravity and causality checks (g1--g3, c1--c2) $43$, the adversarial audit (a1--a5, \cref{app:repro}) $58$ and the release audit
(a6--a8) $20$, of which $9$ are numerical. The complete suite, $61$ scripts, is run by \path{verification/run_all.py}; in the runs of
\path{reproducibility/selected-results/v1.0.0/} the manifest of \cref{app:manifest} lists $1389$ checks, all passing, of which $169$ are numerical, and
every script with an earlier recorded run reproduces its identifiers and outcomes. The runs record the interpreter and library versions. Of the
$1389$ checks, $643$ are cited in the manuscript by identifier; \cref{tab:theoremcode} maps the others to the statements they test.

\paragraph{What a passing check shows.} A \texttt{PASS} line shows that the coded assertion holds: exactly, for identities, ranks and symbolic
derivations, and within the printed tolerance at the sampled inputs, for numerical checks. It validates a statement of the text only to the
extent that the assertion is the statement. A pass does not change the class of a claim: an inherited statement that a script re-derives
remains inherited, and an interpretive reading is not tested at all. The number of \texttt{PASS} lines is not a number of independent claims,
since derivation and referee scripts check many statements twice.

\subsection{What has no executable evidence}

\begin{enumerate}[label=(\roman*),leftmargin=2.2em]
\item \emph{Statements without a cited check.} Of the $438$ theorems, propositions, lemmas and corollaries of the manuscript, the
claim-by-claim map \path{verification/COVERAGE.tsv} lists those that cite a current check, those with a written proof in the text, and those
that rest on the argument given where they are stated, on the review adjudications of \cref{ch:ch13_compatibility}, or on a check that the
text does not yet cite. The algebraic identities and the explicit reductions of \cref{ch:ch11_curvature,ch:ch12_causal} are checked exactly by
the scripts of \path{verification/ch11_curvature_gravity} and \path{verification/ch12_causality} (g1--g3, c1--c2); their gravitational and causal
readings are not, since no gravitational action, interacting causal theory or stability analysis is constructed.
\item \emph{The vacuum analysis.} Since the audit of this version, the second variation on the tensor model, the $271$-dimensional angular
Hessian, the random-start minimisations on $\mathcal T$ and on $\mathcal M$ and the crossing value $x_*$ are reproduced by
\path{verification/ch07_mass_vacuum/} (v1, v2, v3), with archived seeds and tolerances; the audit found and corrected one wrong entry of
\cref{thm:secondvar}. The minimisations remain numerical experiments, and the global statements that rest on them remain conjectures
(\cref{prop:randomstart}). The invariant-theory multiplicities quoted after \cref{thm:potential} are certified exactly by z4-4.
\item \emph{What the release audit leaves.} The electroweak statements of \cref{ch:ch07_vacuum} (\cref{prop:tensor,prop:EW,prop:rankfive,prop:bidoublet}
and the structural part of \cref{thm:mA}) are recomputed numerically (a8), not certified exactly. In \cref{prop:tensor} the identification of the
spectrum of $\Phi^{\!\top}\Phi$ with that of $S^\dagger S$ uses an explicit isomorphism that no check reproduces; the one-loop formula of
\cref{prop:coupled3} is inherited; and \cref{ex:subC} has no check.
\item \emph{Literature.} Bibliographic data were confirmed on publisher, repository or citing pages; where only secondary confirmation was
possible, the author's literature record says so. One attribution, the precise theorem for Howe's skew duality, could
not be checked against the primary source.
\end{enumerate}

\subsection{Corrections found by the checks}

The checks changed the mathematics, not only its presentation. Examples, each recorded with its evidence in the review register: the
surplus-time constraints remove every surplus mode, not only the unstable band (CORE-07); the ordering in the bilinear identity
$Y^{\!\top}TY=-TY^2$ (LAD-01); the binding gap is not a binding energy, and an earlier deficit formula was false (LAD-04); two false statements
about which duplications preserve $|q_2|$; a proposed sign change of the copy-rotation convention was rejected (z1-1d); the factor between
the vector-mass conventions (CORE-01); the counterexample to the vacuum-minimum hypotheses (CORE-02); and the local-stability exception in the
photon family (PH-02). A list of this kind is evidence that the procedure finds errors; it is not evidence that none remain.
